\documentclass[11pt,a4paper]{article}
\usepackage[margin=0.85in,top=1in,bottom=1in]{geometry}
\usepackage{amsmath,amssymb,amsfonts}
\usepackage{graphicx}
\usepackage{float}
\usepackage{enumitem}
\usepackage{microtype}
\usepackage{caption}
\usepackage[hidelinks]{hyperref}
\usepackage[most,breakable]{tcolorbox}
\tcbuselibrary{breakable,skins}
\usepackage[utf8]{inputenc}
\usepackage[T1]{fontenc}
\usepackage{lmodern}
\captionsetup{font=small,labelfont=bf,skip=4pt}
\setlist{leftmargin=1.5em,topsep=2pt}

%% Breakable card — prevents content cutoff across pages
\newtcolorbox{solcard}[3]{
  breakable, enhanced,
  colback=white, colframe=gray!50, arc=2pt,
  title={\textbf{#1}\hfill\textsf{\small #3\,pts}},
  fonttitle=\normalsize\sffamily\bfseries,
  before upper={\small\textit{#2}\par\smallskip\hrule height 0.4pt\smallskip},
  top=4pt, bottom=6pt, left=7pt, right=7pt,
  parbox=false,
}

\title{\textbf{Friction in elastic contact --- Solution}}
\author{\normalsize X25 (2025) — English translation}
\date{}
\begin{document}\maketitle\vspace{-0.5em}
\begin{solcard}{A1}{Give the dimensions of Young's modulus $E$ and Poisson's ratio $\nu$.}{0.10}
\par\smallskip\noindent\textbf{Answer:} \[ [E] = \mathrm{N} / \mathrm{m}^2 = \mathrm{Pa} \quad [\nu] = 1\]
\end{solcard}

\begin{solcard}{A2}{Express the value of $p_0$ in terms of $F$ and $a$.}{0.50}

\[F = \int\limits_0^a p_0 \sqrt{1 - \frac{r^2}{a^2}} 2 \pi r \, dr = \frac{2}{3} \pi p_0 a^2\]

\par\smallskip\noindent\textbf{Answer:} \[p_0 = \frac{3F}{2\pi a^2}\]
\end{solcard}

\begin{solcard}{A3}{Find the moment of dry friction force $M$ relative to the axis of rotation. Express the answer in terms of $F$, $R$, $\mu$, $E$, $\nu$ and the numerical constant $J$ given by \[J=\int\limits_0^1 u^2 \sqrt{1-u^2} du\]}{0.40}

\[\mathrm{d}M =\mu\cdot 2\pi  r^2\cdot p(r)\mathrm{d}r \]\[M = \int\limits_0^a \mu p_0 \sqrt{1-\frac{r^2}{a^2}} 2\pi r^2 \, dr = 2\pi a^3 p_0 \mu J = 3 F a \mu J \]At the same time, from the Hertz formula $a = \sqrt[3]{\frac{3FR(1-\nu^2)}{4E}}$, so in the end

\par\smallskip\noindent\textbf{Answer:} \[M = 3 \sqrt[3]{\frac{3R(1-\nu^2)}{4E}} \mu J F^{4/3} \]
\end{solcard}

\begin{solcard}{A4}{Obtain an explicit form of the function $J(x)$ given by \[J(x) = \int\limits_0^x u^2 \sqrt{1-u^2} du.\] Note: if you are unable to find an explicit form $J(x)$, then you can use $J(1)=\pi/16$.}{0.50}

Note: If you are unable to find the explicit form $J(x)$, then you can use $J(1)=\pi/16$.


Using the replacement $u= \sin v$ ($du =  \cos v \, dv$) we get \[J(x) = \int\limits_0^{\arcsin x} \sin^2 v \cos^2 v \, dv = \frac{1}{4} \int\limits_0^{\arcsin x} \sin^2 2v \, dv = \frac{1}{8} \int\limits_0^{\arcsin x} \left( 1 - \cos 4 v \right) dv = \frac{\arcsin x}{8} - \frac{\sin ( 4 \arcsin x )}{32} \]In this case, \[\sin 4x = 2 \sin 2x \cos 2x = 4 \sin x \cos x (1-2 \sin^2 x) = 4 \sin x \sqrt{1-\sin^2 x} (1- 2 \sin^2 x),\] means \[J(x) = \frac{\arcsin x}{8} -\frac{x\sqrt{1-x^2}(1-2x^2)}{8}\]

\par\smallskip\noindent\textbf{Answer:} \[ J(x) =\frac{\arcsin x}{8} -\frac{x\sqrt{1-x^2}(1-2x^2)}{8}\]
\end{solcard}

\begin{solcard}{B1}{Using an oscilloscope, measure the sequence of times $t_i$ passing the black stripe over the photodiode. There must be at least $15$ times in the sequence!}{0.50}

$i$ 1 2 3 4 5 6 7 8 9 10 11 12 13 14 15 $t_i, \mathrm{s}$ 0.10 0.22 0.34 0.45 0.56 0.68 0.78 0.91 1.03 1.15 1.28 1.41 1.53 1.66 1.79 $i/t_i, \mathrm{s}^{-1}$ 9.62 9.26 8.93 8.93 8.93 8.82 8.93 8.77 8.74 8.70 8.59 8.51 8.50 8.43 8.38

\end{solcard}

\begin{solcard}{B2}{Construct a graph of $t_i$ versus $i$ in such coordinates that it is linear. Find the parameters of the line.}{0.50}
\begin{figure}[H]
  \centering
  \includegraphics[width=0.72\textwidth]{E3_S_EN_assets/fig1.jpg}
\end{figure}
\end{solcard}

\begin{solcard}{B3}{On the answer sheet, indicate by ticking which of the indicated combinations will cause the rotation of the top to be stable.}{0.30}
\end{solcard}

\begin{solcard}{B4}{For 6 different stable combinations of nuts and bolts bolted to the attachment, take measurements similar to step B1. Draw the combinations used in the way suggested in the condition. Note: measurements are completely similar to part B1, i.e. an attachment without bolts and nuts does NOT count as a combination. Combinations that do not differ in terms of dynamics are considered one.}{3.00}

Note: measurements are completely similar to part B1, i.e. an attachment without bolts and nuts does NOT count as a combination. Combinations that do not differ in terms of dynamics are considered one.


Combination 1


Combination 2


Combination 3


Combination 4


Combination 5


Combination 6

\par\smallskip\noindent\textbf{Answer:} $i$ 1 2 3 4 5 6 7 8 9 10 11 12 13 14 15 $t_i, \mathrm{s}$ 0.084 0.168 0.256 0.344 0.432 0.524 0.612 0.708 0.800 0.896 0.992 1.096 1.188 1.296 1.396 $i/t_i, \mathrm{s}^{-1}$ 11.905 11.905 11.719 11.236 11.261 11.194 11.218 11.111 11.084 11.013 10.956 10.830 10.833 10.703 10.653
\begin{figure}[H]
  \centering
  \includegraphics[width=0.72\textwidth]{E3_S_EN_assets/fig2.jpg}
\end{figure}
\begin{figure}[H]
  \centering
  \includegraphics[width=0.72\textwidth]{E3_S_EN_assets/fig3.jpg}
\end{figure}
\begin{figure}[H]
  \centering
  \includegraphics[width=0.72\textwidth]{E3_S_EN_assets/fig4.jpg}
\end{figure}
\begin{figure}[H]
  \centering
  \includegraphics[width=0.72\textwidth]{E3_S_EN_assets/fig5.jpg}
\end{figure}
\begin{figure}[H]
  \centering
  \includegraphics[width=0.72\textwidth]{E3_S_EN_assets/fig6.jpg}
\end{figure}
\begin{figure}[H]
  \centering
  \includegraphics[width=0.72\textwidth]{E3_S_EN_assets/fig7.jpg}
\end{figure}
\end{solcard}

\begin{solcard}{B5}{Using the least squares method, conduct an analysis of the dependencies obtained in the previous paragraph, similar to your solution to paragraph B2.}{0.30}
\par\smallskip\noindent\textbf{Answer:} Combination 1
\begin{figure}[H]
  \centering
  \includegraphics[width=0.72\textwidth]{E3_S_EN_assets/fig8.jpg}
\end{figure}
\begin{figure}[H]
  \centering
  \includegraphics[width=0.72\textwidth]{E3_S_EN_assets/fig9.jpg}
\end{figure}
\begin{figure}[H]
  \centering
  \includegraphics[width=0.72\textwidth]{E3_S_EN_assets/fig10.jpg}
\end{figure}
\begin{figure}[H]
  \centering
  \includegraphics[width=0.72\textwidth]{E3_S_EN_assets/fig11.jpg}
\end{figure}
\begin{figure}[H]
  \centering
  \includegraphics[width=0.72\textwidth]{E3_S_EN_assets/fig12.jpg}
\end{figure}
\begin{figure}[H]
  \centering
  \includegraphics[width=0.72\textwidth]{E3_S_EN_assets/fig13.jpg}
\end{figure}
\end{solcard}

\begin{solcard}{B6}{Based on the results from B2 and B5, construct a linearized graph to find the Young's modulus $E$ of polybutadiene rubber.}{0.40}

The angular acceleration $\dot{\omega}$ is related to the moment of inertia $I$ and the frictional moment $M$ through the basic rotational dynamics equation: \[ I \dot{\omega} = M\]The moment of inertia $I$ and the total mass $m$ can be calculated for each bolt and nut combination. Then the coordinates that linearize the dependence are $I \dot{\omega}$ of $m^{4/3}$. The straight line must pass through zero and have a slope \[k= 3 \sqrt[3]{\frac{3R(1-\nu^2)}{4E}} \mu J g^{4/3}\]

\begin{figure}[H]
  \centering
  \includegraphics[width=0.72\textwidth]{E3_S_EN_assets/fig14.jpg}
\end{figure}
\end{solcard}

\begin{solcard}{B7}{Find the Young's modulus $E$ of polybutadiene rubber.}{0.50}
\par\smallskip\noindent\textbf{Answer:} \[E = 108.6~\mathrm{MPa}\]
\end{solcard}

\end{document}